\documentclass[10pt,a4paper]{amsart}
\usepackage[margin=19mm]{geometry}
\usepackage{amsmath,amssymb,mathtools}
\usepackage[scr=rsfs,cal=euler]{mathalfa}
\usepackage{xfrac}
\usepackage[unicode,hidelinks,bookmarks=false]{hyperref}
\newcommand{\ie}{i.e.\ }
\newcommand{\cf}{cf.\ }
\newcommand{\resp}{respectively\ }
\newcommand{\Subsection}{\subsection}
\providecommand{\nobreakdash}{\nobreak\hbox{-}}
\setlength{\emergencystretch}{2em}
\multlinegap0pt
\usepackage{bm}
\usepackage{bbm}
\usepackage{dsfont}
\usepackage{xspace}
\usepackage{cancel}
\usepackage{mathtools}
\usepackage{amsthm}
\makeatletter
\@ifundefined{theorem}{\newtheorem{theorem}{Theorem}}{}
\@ifundefined{lemma}{\newtheorem{lemma}{Lemma}}{}
\@ifundefined{proposition}{\newtheorem{proposition}{Proposition}}{}
\makeatother
\newcommand{\D}{\mathbb{D} }
\newcommand{\F}{\mathbb{F} }
\newcommand{\R}{\mathbb{R} }
\newcommand{\Z}{\mathbb{Z} }
\newcommand{\bG}{\mathbf{G}}
\newcommand{\bS}{\mathbf{S}}
\newcommand{\fraka}{\mathfrak{a}}
\newcommand{\frakg}{\mathfrak{g}}
\newcommand{\tran}{{^{\mathsf{T}}\!}}
\title[Semisimple algebraic groups over real closed fields]{Semisimple algebraic groups over real closed fields}
\author{Raphael Appenzeller}
\date{Source: arXiv:2601.07732v1 (CC BY 4.0)}
\begin{document}
\maketitle

\noindent\emph{Source and licence.} Semisimple algebraic groups over real closed fields. Version arXiv:2601.07732v1 (CC BY 4.0), distributed under the Creative Commons Attribution 4.0 International licence, as recorded in the arXiv licence metadata for this exact version and shown on the abstract page https://arxiv.org/abs/2601.07732v1. Full rights notice: licenses/arxiv-2601.07732-CC-BY-4.0.txt.\par\smallskip

\noindent\textit{Editorial scope.} Counted unit: the Lemma whose source label is \texttt{lem:Jacobson\_Morozov\_m} in the article (source0.tex lines 1825--1842, source label \texttt{lem:Jacobson\_Morozov\_m}), delivered together with the article's own complete original proof of it (source0.tex lines 1843--1917, one \texttt{proof} environment of 75 lines). No mathematics is written, repaired or re-derived: statement and proof are the article's own text line for line. Required context retained uncounted: lem:JM\_basic (lines 864--869); lem:char (lines 1302--1309); lem:aexpXa (lines 1559--1563); prop:Jacobson\_Morozov\_real\_closed (lines 1732--1750); lem:Jacobson\_Morozov\_oneparam (lines 1787--1794). Every definition, display and label of those lines is retained unchanged. Omitted from this package and not redistributed: 1--863, 870--1301, 1310--1558, 1564--1731, 1751--1786, 1795--1824, 1918--2218. Nothing inside a retained excerpt is omitted or altered: every retained body is delivered exactly as the article's lines are written, so no omission receipt of any kind accompanies this package. Citations: the retained excerpts carry no citation command at all, so no bibliography entry of the article is transcribed and no reference is redistributed. Reference adaptation: none was needed and none was made; every reference inside the delivered unit resolves inside it, so no unresolved reference or citation is produced. Printed slips kept literal: no printed word, symbol, hypothesis or number of the counted statement or of its proof was altered. Layout and class: the article's own class is \texttt{amsart} with packages including geometry, amsmath, amscd, amssymb, xcolor, pst-node, tikz-cd, graphicx, comment, parskip, hyperref, fancyhdr, subcaption, amsthm; the wrapper is the lane's pinned \texttt{amsart} editorial wrapper at 10pt on A4, which restates the theorem-like environments the retained text needs and adds the article's own macro definitions that the retained text needs (\texttt{\textbackslash D}, \texttt{\textbackslash F}, \texttt{\textbackslash R}, \texttt{\textbackslash Z}, \texttt{\textbackslash bG}, \texttt{\textbackslash bS}, \texttt{\textbackslash fraka}, \texttt{\textbackslash frakg}, \texttt{\textbackslash tran}), with extra wrapper packages (bm, bbm, dsfont, xspace, cancel, mathtools). The delivered numbering is the unit's own. Assets: no retained span contains a figure environment, a graphics-inclusion command, a PDF-page inclusion, a table or a TikZ picture; no image file is copied into this package, the package declares no asset dependency and no image file is redistributed with it. Licence: version arXiv:2601.07732v1 of this article is distributed under the Creative Commons Attribution 4.0 International licence, as recorded in the arXiv licence metadata for this exact version and shown on the abstract page https://arxiv.org/abs/2601.07732v1. A full rights notice is delivered at licenses/arxiv-2601.07732-CC-BY-4.0.txt. Characters: every retained excerpt is pure ASCII, so the delivered engine, XeLaTeX with its default Latin Modern text font, reports no missing character.
\begin{lemma}\label{lem:JM_basic}
	Let $\alpha \in \Sigma$ and $X \in \frakg_\alpha \setminus \{0\}$. Then there is a $Y \in \frakg_{-\alpha} \setminus \{0\}$ and $H \in \fraka$ such that $(X,Y,H)$ forms a $\mathfrak{sl}(2)$-triplet, i.e.
	$$
	[X,Y] = H, \quad [H,X] = 2X,  \quad [H,Y]=2Y.
	$$
\end{lemma}

\begin{lemma}\label{lem:char}
	Let $\alpha \in \Sigma$. Then the homomorphism
	\begin{align*}
		\chi_\alpha \colon A_\R & \to \R_{>0} \\
		\exp(H) &\mapsto e^{\alpha(H)}
	\end{align*}
	is the restriction to $A_\R$ of an algebraic character in $\hat{\bS}$. 
\end{lemma}

\begin{lemma} \label{lem:aexpXa}
	Let $\alpha \in \Sigma, X\in \mathfrak{g}_\alpha, X' \in \mathfrak{g}_{2\alpha}$ and $a \in A_\F $. Then 
	$$a\exp\left(X + X'\right)a^{-1} = \exp\left(\chi_\alpha (a) X + \chi_\alpha(a)^2 X'\right),$$
	where $\chi_\alpha \colon A_\F \to \mathbb{F}_{>0}$ is the algebraic character from Lemma \ref{lem:char}.
\end{lemma}

\begin{proposition}\label{prop:Jacobson_Morozov_real_closed}
	Let $\alpha \in \Sigma$ and assume $\frakg_{2\alpha} = 0$. Let $u \in (U_{\alpha})_\F$. Then there are $ X \in (\frakg_\alpha)_\F$ and $t \in \F$ such that $u=\exp(tX)$. Let $(X,Y,H)$ be the $\mathfrak{sl}_2$-triplet of Lemma \ref{lem:JM_basic}. Then there is a group homomorphism $\varphi_\F \colon \operatorname{SL}(2,\F) \to G_\F$ that is a restriction of a morphism of algebraic groups
	$
	\varphi \colon \operatorname{SL}(2,\D)  \to \bG 
	$
	defined over $\F$ such that $\varphi$ has finite kernel and
	\begin{align*}
		\varphi_\F \begin{pmatrix}
			1 & t \\ 0& 1
		\end{pmatrix} =  u = \exp(tX)  \quad \text{and} \quad  \varphi_\F \begin{pmatrix}
		1 & 0 \\ t & 1
		\end{pmatrix} = \exp(tY).
	\end{align*}
	If $\varphi$ is not injective, then $\ker (\varphi) \cong \Z/2\Z$ and $\varphi$ factors through the isomorphism
	$$
	\operatorname{PGL}(2,\D) := \operatorname{SL(2,\D)}/\ker(\varphi)  \xrightarrow{\sim} \varphi(\operatorname{SL}(2,\D))
	$$
	which is also defined over $\F$. Moreover $\varphi(g\tran) = \varphi(g)\tran$, for any $g \in \operatorname{SL}(2,\F)$.
\end{proposition}

\begin{lemma}\label{lem:Jacobson_Morozov_oneparam}
	Let $\alpha \in \Sigma$ such that $(\frakg_{2\alpha})_\F = 0$. Let $u\in (U_\alpha)_\F$ and $\varphi \colon \operatorname{SL}(2,\D) \to \bG$ as in Proposition \ref{prop:Jacobson_Morozov_real_closed}. For every $\lambda \in \F_{>0}$
	$$
	\chi_\alpha \left(\varphi_\F \begin{pmatrix}
		\lambda & 0 \\ 0 & \lambda^{-1}
	\end{pmatrix} \right) = \lambda^2.
	$$
\end{lemma}

\begin{lemma}\label{lem:Jacobson_Morozov_m}
	Let $\alpha \in \Sigma$ such that $(\mathfrak{g}_{2\alpha})_\F = 0$. Let $u \in (U_\alpha)_\F $ and $X\in (\frakg_\alpha)_\F$, $t\in \F$ and $\varphi \colon \operatorname{SL}(2,\D) \to \bG$ as in Proposition \ref{prop:Jacobson_Morozov_real_closed}, so that $u = \exp(tX)$. The element
	$$
	m(u) := \varphi_\F \begin{pmatrix}
		0 & t \\ -1/t & 0
	\end{pmatrix} \in G_\F
	$$ 
	is contained in $\operatorname{Nor}_{G_\F}(A_\F)$ and 
	$$
	\chi_\alpha(m(u)\cdot a\cdot m(u)^{-1}) = \chi_\alpha(a)^{-1}
	$$ 
	for any $a \in A_\F$.
	
	When $t=1$, even $m(u) \in N_\F = \operatorname{Nor}_{K_\F}(A_\F)$ and 
	$$
	m(u) \cdot (U_\alpha)_\F \cdot  m(u)^{-1} = (U_{-\alpha})_\F.
	$$
\end{lemma}

\begin{proof}
	Let $a\in A_\F$ and
	$$
	a_\perp := \varphi_\F \begin{pmatrix}
		\sqrt{\chi_\alpha(a)} & 0 \\ 0 & \sqrt{\chi_\alpha(a)}^{-1}
	\end{pmatrix},
	$$ 
	then $a_0:= aa_\perp^{-1}$ satisfies
	$$
	\chi_\alpha (a_0) = \chi_\alpha(a) \chi_\alpha(a_\perp)^{-1} = \chi_\alpha(a) \frac{1}{\chi_\alpha(a)} = 1 
	$$
	by Lemma \ref{lem:Jacobson_Morozov_oneparam}. For any $u' \in \exp((\frakg_{\alpha})_\F)$ or $u'\in \exp((\frakg_{-\alpha})_\F)$, $a_0u'=u'a_0$ by Lemma \ref{lem:aexpXa}. We note that
	\begin{align*}
			m(u) &= \varphi_\F \begin{pmatrix}
			0 & t \\ -1/t & 0
		\end{pmatrix} =  \varphi_\F \left( \begin{pmatrix}
			1 & 0 \\ -1/t & 1
		\end{pmatrix}  \begin{pmatrix}
			1 & t \\ 0 & 1
		\end{pmatrix} \begin{pmatrix}
			1 & 0 \\ -1/t & 1
		\end{pmatrix}\right) \\
		& \in \exp((\frakg_{-\alpha})_\F) \cdot u \cdot \exp((\frakg_{-\alpha})_\F)
	\end{align*}
	which implies $a_0m(u) = m(u)a_0$. Now
	\begin{align*}
		m(u) \cdot a \cdot m(u)^{-1} &= m(u) \cdot a_0 \cdot a_\perp \cdot m(u)^{-1} = a_0 \cdot m(u) \cdot  a_\perp \cdot m(u)^{-1} \\
		&= a_0 \cdot \varphi_\F \left(\begin{pmatrix}
			0 & t \\ -1/t & 0
		\end{pmatrix}\begin{pmatrix}
		\sqrt{\chi_\alpha(a)} & 0 \\ 0 & \sqrt{\chi_\alpha(a)}^{-1}
		\end{pmatrix}\begin{pmatrix}
		0 & -t \\ 1/t & 0
		\end{pmatrix}\right) \\
		&= a_0\cdot \varphi_\F \begin{pmatrix}
		\sqrt{\chi_\alpha(a)}^{-1} & 0 \\ 0 & \sqrt{\chi_\alpha(a)}
		\end{pmatrix} = a_0 a_\perp^{-1} \in A_\F
	\end{align*}
	and thus $m(u) = \operatorname{Nor}_{G_\F}(A_\F)$. We see directly that 
	$$
	\chi_\alpha (m(u) \cdot a\cdot m(u)^{-1}) = \chi_\alpha(a_0 \cdot a_\perp^{-1}) = \chi_\alpha(a_0) \cdot \chi_\alpha(a_\perp)^{-1} = \chi_\alpha(a)^{-1}.
	$$
	Now if $t=1$ we can use that $\varphi$ preserves transposition by Proposition \ref{prop:Jacobson_Morozov_real_closed}, to show
	\begin{align*}
		m(u) \cdot m(u)\tran &= \varphi_\F \begin{pmatrix}
			0 & 1 \\ -1 & 0
		\end{pmatrix} \left( \varphi_\F \begin{pmatrix}
			0 & 1 \\ -1 & 0
		\end{pmatrix}\right)\tran \\
		& = \varphi_\F \left( \begin{pmatrix}
			0 & 1 \\ -1 & 0
		\end{pmatrix}   \begin{pmatrix}
			0 & -1 \\ 1 & 0
		\end{pmatrix} \right)= \varphi_\F\begin{pmatrix}
			1 & 0 \\ 0 & 1
		\end{pmatrix} = \operatorname{Id}
	\end{align*}
	and hence $m(u)\in K_\F$, so $m(u) \in N_{\F} = \operatorname{Nor}_{K_\F}(A_\F)$. Thus $m(u)$ is a representative of an element $w=[m(u)]$ of the spherical Weyl group $W_s = N_\F/M_\F$. While $u\in (U_{\alpha})_\F$,
	\begin{align*}
		m(u)\cdot u\cdot m(u)^{-1} &= \varphi_\F \left( \begin{pmatrix}
			0 & 1 \\ -1 & 0
		\end{pmatrix}\begin{pmatrix}
			1 & t \\ 0 & 1
		\end{pmatrix} \begin{pmatrix}
			0 & -1 \\ 1 & 0
		\end{pmatrix} \right) \\
		&= \varphi_\F \begin{pmatrix}
			1 & 0 \\ -t & 1
		\end{pmatrix} \in (U_{-\alpha})_\F,
	\end{align*}
	and thus we have $w(\alpha) = -\alpha$. Then 
$$
m(u) \cdot (U_\alpha)_\F \cdot  m(u)^{-1} = (U_{w(\alpha)})_\F =  (U_{-\alpha})_\F.
$$
\end{proof}

\end{document}
